% ===========================================================================
\section{Implementation Plan}
% Source material: docs/roadmap.md, docs/architecture.md
% ===========================================================================

We will build tier 2 first and then tier 3 (Table~\ref{tab:tiers}). Tier 3
replaces tier 2's Merkle tree, while the phases, the factory, the indexer and
the gas baseline all carry over.

The reason for building tier 2 first is that the cost of anonymity is a
difference between the two tiers. Tier 3 on its own would tell us what a ballot
costs, but not how much of that cost comes from anonymity. The gas snapshot
taken at the end of Phase 1 gives one number, the one from Phase 2 gives the
other, and A1 reports the difference. C1 then plots both tiers: tier 3 should
stay flat where tier 2 rises, and cost more throughout. Tier 2 has no proving
step, so B2 and C2 measure tier 3 only.

\subsection{Phases and milestones}

Implementation starts after the mid-semester examinations. The work is split
into five phases, done in order, with everything finished before Evaluation~2
(Table~\ref{tab:phases}).

\begin{table}[htbp]
  \centering
  \caption{Phase plan and milestones.}
  \label{tab:phases}
  \small
  \begin{tabular}{@{}lL{14.0cm}@{}}
    \toprule
    \textbf{Phase} & \textbf{Milestone}                                                        \\
    \midrule
    1              & M1: a vote works end to end locally (tier 2, without anonymity or gas sponsorship) \\
    2              & M2: anonymous and gasless voting on the testnet (tier 3)                  \\
    3              & Feature freeze; auditing, resilience and usability testing                \\
    4              & Measurements for every parameter in Section 5                             \\
    5              & Demo rehearsal and buffer time                                            \\
    \bottomrule
  \end{tabular}
\end{table}

\begin{itemize}
  \item \emph{Phase 1} delivers milestone M1, a complete vote from start to
        finish. We do this first because projects that leave integration to
        the end usually run into trouble. To pass, we must be able to start
        from a fresh copy of the code, deploy to a local node (Anvil), register
        a few voters, cast their ballots from the browser and read back the
        tally, with the contracts covered by tests. The first gas snapshot is
        recorded the same day.
  \item \emph{Phase 2} tackles the two hardest parts early: the proofs and the
        relayer, including the check that rejects ballots beyond the voter
        count. M2 is reached when a voter with no ETH can cast an anonymous
        ballot on Sepolia, and nothing on-chain links their address to it.
  \item \emph{Phase 3} is about finishing rather than adding features. It
        covers \texttt{TallyLib} with plurality voting (and approval voting
        only if time allows), a clean report from the Slither static analyser,
        handling of chain reorganisations, auditing with the backend switched
        off, and the E3 usability test. After this feature freeze we only fix
        bugs, run benchmarks and write.
  \item \emph{Phase 4} measures every parameter in Section~5. Each parameter
        ends up with a number, or with a clear reason why it has none.
  \item \emph{Phase 5} is for rehearsing the demo on Sepolia. During the demo
        we register everyone in the room and cast a batch of ballots before any
        live one, so that no ballot is the only one in its time window.
\end{itemize}

\subsection{Risks and contingencies}

Our biggest risk is generating proofs in the browser, and on a phone this is
more likely to fail completely than to be merely slow. Proving runs in a
background thread, and a device that cannot finish a proof counts as a failure
of B2.

We will make a go/no-go decision one week after reaching M1. If phones still
cannot produce valid proofs by then, proving moves to a server. This fallback
is weaker, and we would report it as such: the voter's secret would leave their
device, so the server could link a voter to their ballot. Anonymity would still
hold against anyone reading the blockchain, but not against our own operator.
Table~\ref{tab:risks} lists the other main risks and how we would respond.

\begin{table}[htbp]
  \centering
  \caption{Risks and responses.}
  \label{tab:risks}
  \small
  \begin{tabular}{@{}L{4.0cm}L{9.8cm}@{}}
    \toprule
    \textbf{Risk}            & \textbf{Response}                                                                                                             \\
    \midrule
    Tier 3 not ready in time & Decide a week before the feature freeze. If tier 3 is not ready, demonstrate tier 2 and report tier 3 as unfinished. Tier 2 is a complete system, so the worst case is a demo without anonymity \\
    Relayer down on the day  & Voters submit directly from funded wallets; this affects E2, not the election itself                                              \\
    Relayer funds drained    & Before paying, the relayer checks the proof and that the nullifier is unused, so it pays for at most one ballot per registered voter \\
    Not enough Sepolia ETH   & Faucets cannot fund heavy load, so load and gas tests run on a local node, and Sepolia is used only for latency tests and the demo \\
    \bottomrule
  \end{tabular}
\end{table}
